% ECONOMICS_PROBLEM: {"id": "D31-379", "title": "Inheritance and long-run inequality", "jel_code": "D31", "source_id": "OP-0369"}
% !TeX program = xelatex
\documentclass[11pt,a4paper]{article}
\usepackage[margin=23mm]{geometry}
\usepackage{fontspec}
\setmainfont{Latin Modern Roman}
\usepackage{amsmath,amssymb,mathtools}
\usepackage{xcolor,enumitem,booktabs,longtable,array,xurl,hyperref}
\definecolor{muted}{HTML}{53636C}
\hypersetup{colorlinks=true,urlcolor=blue,linkcolor=blue}
\setlength{\parindent}{0pt}
\setlength{\parskip}{5pt}
\newcommand{\R}{\mathbb R}
\newcommand{\E}{\mathbb E}
\newcommand{\N}{\mathbb N}
\newcommand{\ind}{\mathbf 1}
\newcommand{\Var}{\operatorname{Var}}
\newcommand{\Cov}{\operatorname{Cov}}
\newcommand{\supp}{\operatorname{supp}}
\DeclareMathOperator*{\argmin}{arg\,min}
\DeclareMathOperator*{\argmax}{arg\,max}
\newcommand{\entrymeta}[3]{{\sffamily\small\color{muted}\textbf{\texttt{#1}}\quad|\quad #2\par #3\par}}
\newcommand{\Problemref}[1]{\texttt{#1}}
\newcommand{\JELref}[2]{\texttt{#2} (JEL \texttt{#1})}
\begin{document}
\section*{Inheritance and long-run inequality}
\textbf{Problem ID:} \texttt{D31-379}\quad \textbf{JEL:} \texttt{D31}\par
% BEGIN ECONOMICS STATEMENT
\label{problem:OP-0369}
{\sffamily\small\color{muted}Original subject group: Inequality and economic mobility\par}
\label{definitions:problem:30007651}
\textbf{Audit classification:} Background; proposed extension. Authors, 2025 year, title and DOI checked. The study estimates marginal transfer effects; dynamic fiscal counterfactuals are editorial, not an explicitly verified remaining conjecture.
\subsection*{Definitions and precise research target}
\textbf{Complete editorial problem.} Fix an adult cohort with pre-inheritance wealth, parental links, and longitudinal portfolio records. Let \(b_i\) be person \(i\)'s net inheritance under a declared timing rule, and let \(b=(b_1,\ldots,b_N)\) be a feasible population-wide assignment for the \(N\) adults. Let \(W_{it}(b)\) be subsequent wealth allowing interactions among recipients. Its evolution includes labor income, consumption, realized portfolio returns, gifts, and mortality. For the population wealth distribution \(F_t(b)\), define \(\Delta_I(t,b)=I(F_t(b))-I(F_t(0))\), where \(I\) is a specified relative inequality measure. Specify the resource-feasible no-inheritance comparator, including who receives estates instead. Estimate this path for different inheritance sizes, prior wealth ranks, and recipient ages. Determine how much variation in the path is accounted for by heirs' consumption and saving responses, business investment, portfolio choice, and returns, using additional identifying variation rather than accounting residuals alone. Preserve any established immediate equalizing or disequalizing effects in the relevant data. The research extension concerns dynamic behavioral mechanisms and their transportability; an observed cross-sectional marginal-transfer effect does not identify the long-run response. Death-related shocks may simultaneously alter grief, work, and transfers, so causal designs must address their exclusion restrictions. The cited 2025 paper answers size-dependent cross-sectional questions but an explicit remaining open-question passage for this dynamic formulation was not verified.

\textbf{Definitions and admissible scope.} An inheritance intervention specifies donors as well as recipients, mortality timing, estate taxation, inter vivos gifts, and the use of revenues. The vector \(b=0\) must not silently destroy resources: choose a no-inheritance comparator with a stated recipient or public disposition of estates, or define a mechanical recipient-accounting subtraction and label it descriptive. Wealth \(W_t(b)\) includes borrowing and all behavioral responses, and the same baseline adults are followed with an explicit treatment of deaths. Choose an inequality functional whose domain accommodates negative wealth, such as a quantile-integral top share with finite absolute wealth and positive mean; a conventional nonnegative Gini needs stronger restrictions. A death instrument must rule out the donor's death affecting outcomes other than through the targeted transfer if that exclusion is used. Take \(N<\infty\), a fixed baseline-adult cohort, a finite horizon \(t\), and positive normalized weights. Specify how deceased members' wealth is assigned for population comparison. For example choose \(I(F)=\int_{1-p}^1F^{-1}(u)\,du/E[W]\), with \(p\in(0,1)\), finite absolute wealth, and positive mean under all policies. The baseline symbol \(F_t(0)\) means a complete alternative estate-disposition regime giving no inheritance to the focal recipients; its displaced estates go to the declared other recipients or public budget. A purely mechanical removal holding behavior fixed is an explicitly different exercise. The dynamic \(\Delta_I\) requires an admissible causal law matching inheritance timing and wealth histories, not only the observed cross-sectional transfer distribution.
\subsection*{Variants and assumption changes}
\begin{enumerate}
\item Cross-sectional marginal transfer (source-based scope): study small changes in the distribution of observed transfer recipients and transfer sizes. The cited paper supplies influence-function estimates for this exercise; it is an existing answer under its assumptions.
\item Dynamic fiscal counterfactual (editorial): vary estate-tax exemption and specify how revenue is transferred to other adults. Estimate saving, investment, and equilibrium prices over a long horizon, comparing the full wealth distribution rather than a static subtraction.
\end{enumerate}
\subsection*{References and source audit}
\textbf{Support and access.} Authors, 2025 year, title and DOI checked. The study estimates marginal transfer effects; dynamic fiscal counterfactuals are editorial, not an explicitly verified remaining conjecture. \textbf{Locator:} Author institutional abstract and citation, JPubE 247, 105398.
\begin{enumerate}\raggedright
\item\hypertarget{definitions:source:30007651:1}{} Salvatore Morelli; Brian Nolan; Juan C. Palomino; Philippe Van Kerm. 2025. \emph{The influence of inheritances on wealth inequality in rich countries}. Abstract, paragraphs 1–3.
\url{https://www.inet.ox.ac.uk/publications/the-influence-of-inheritances-on-wealth-inequality-in-rich-countries}.
DOI: \url{https://doi.org/10.1016/j.jpubeco.2025.105398}.
\end{enumerate}

\subsection*{Common conventions: Source mathematical conventions}

These conventions apply to each entry unless explicitly replaced.

\textbf{Models and identification.} A primitive model \(m\) specifies all probability laws, feasible choices, information, equilibrium and measurement rules used by its target. The admissible class \(\mathcal M\) is fixed independently of outcomes. If \(P_O(m)\) is its observable probability law and \(\tau(m)\) the target, its identified set at observable law \(P\) is
\[
 \mathcal I_\tau(P)=\{\tau(m):m\in\mathcal M,\ P_O(m)=P\}.
\]
Point identification means that this set is a singleton. An empty set signals incompatibility of data and assumptions. Coordinate bounds need not describe the joint set. Bounds are sharp only if every retained target is attainable or arbitrarily approachable by compatible models. Finite bounds require explicit restrictions on unobserved quantities; unrestricted confounding, hidden wealth, extrapolation or arbitrary equilibrium selection can yield uninformative sets.

\textbf{Causal interventions.} A potential outcome \(Y_i(p)\) is the outcome under a complete policy regime \(p\), including timing, financing, information and market spillovers. The target population and units are fixed before treatment. With interference, use the complete assignment vector or a justified exposure mapping; individual no-interference notation alone cannot identify an economy-wide intervention. A difference of mean potential outcomes does not require the cross-world joint distribution. Conditioning on post-treatment migration, survival, enrollment or employment changes the estimand and needs separate justification.

\textbf{Statistical statements.} An estimator, confidence set, forecast or test specifies a sampling law, model class and loss/coverage criterion. Uniform \((1-\alpha)\)-coverage means \(\inf_{m\in\mathcal M}\Pr_m\{\tau(m)\in C_n\}\geq1-\alpha\), or an explicitly asymptotic version. The whole parameter space trivially has coverage, so informativeness requires a stated width, risk or power objective. A forecasting improvement is relative to a named benchmark, information set and evaluation population.

\textbf{Equilibrium and optimization.} An equilibrium comprises feasible plans, optimality, beliefs where needed, budget constraints and all market/resource clearing. The primitives or a stated benchmark define each correspondence \(\mathcal E(p;m)\). Nonemptiness is assumed or proved, never inferred just from compact policy sets. A selected-equilibrium target uses a declared selection rule; a robust target quantifies over all permitted equilibria. Use suprema or infima unless attainment is established. An \(\varepsilon\)-optimal policy is feasible and within \(\varepsilon\) of the finite optimum value. Report an empty feasible set separately.

\textbf{Welfare and accounting.} Normative utility, interpersonal normalization, social weights, numeraire, discounting and the use of public receipts are primitives. Behavioral utility need not coincide with the welfare criterion. Household welfare includes ownership income and rents once; adding profits again to compensating variation generally double-counts them. A monetary equivalent is defined by an explicit transfer and price convention, with monotonicity and range conditions. Average effects, mechanism contributions, transfers and welfare losses are different targets. Infinite sums require convergence and dynamic feasible plans require terminal or no-Ponzi conditions.

\textbf{Source versions.} A working-paper date, revision date and journal publication year are distinguished. A DOI for a journal version is not automatically the DOI of an earlier manuscript. Locators refer to the stated version and printed pages where specified. A metadata-only check does not verify a claim about a full-text conclusion. Every entry's source subsection narrows the provenance claim accordingly.
% END ECONOMICS STATEMENT
\end{document}
